\documentclass[11pt,reqno]{amsart}
\usepackage{amsmath}
\usepackage{amssymb}
\usepackage{amsthm}
\usepackage{enumerate}
\usepackage[notcite,notref]{showkeys}
\usepackage{kotex}

\usepackage{url}
\usepackage{hyperref}

% PAGE SETTING ---------------------------------------

\textheight 22.5  true cm
\textwidth 15 true cm
\voffset -1.0 true cm
\hoffset -1.0 true cm
\marginparwidth= 2 true cm
\renewcommand{\baselinestretch}{1.2}

% THEOREMS -------------------------------------------

\renewcommand{\thefootnote}{\fnsymbol{footnote}}
\newtheorem{thm}{Theorem}[section]
\newtheorem{cor}[thm]{Corollary}
\newtheorem{lem}[thm]{Lemma}
\newtheorem{prop}[thm]{Proposition}
\newtheorem{slem}[thm]{Sublemma}
\theoremstyle{definition}
\newtheorem{defn}[thm]{Definition}
\newtheorem{step}{Step}
\theoremstyle{remark}
\newtheorem{rem}[thm]{Remark}
\newtheorem{prob}[thm]{Problem}
\numberwithin{equation}{section}

\newtheorem{innercustomthm}{Theorem}
\newenvironment{customthm}[1]
{\renewcommand\theinnercustomthm{#1}\innercustomthm}
{\endinnercustomthm}


\newcommand{\ep}{\epsilon}
\newcommand{\lx}{{\lambda, \xi}}
\newcommand{\pa}{\partial}

\newcommand{\mc}{\mathcal{C}}
\newcommand{\mh}{\mathcal{H}}
\newcommand{\mt}{\mathcal{T}}
\newcommand{\mv}{\mathcal{V}}
\newcommand{\mn}{\mathbb{N}}
\newcommand{\mr}{\mathbb{R}}
\newcommand{\G}{\mathcal{G}}
\newcommand{\E}{\mathcal{E}}
\newcommand{\cW}{\mathcal{W}}

\newcommand{\diag}{{\rm diag}}

\renewcommand{\(}{\left(}
\renewcommand{\)}{\right)}

% MATH ------------------------------------------------

\newcommand{\supp}{\text{supp }}
\newcommand{\inp}[2]{\langle #1,#2 \rangle}
\newcommand{\dist}{\text{dist }}
\newcommand{\ty}{{\widetilde y}}
\newcommand{\R}{\mathbb{R}}

%-----------------------------------------------------------------

\newcommand {\com}[1]{\texttt{($\leftarrow$ #1)}}
\newcommand{\inner}[2]{\left< #1, #2 \right>}

\newenvironment{text-box}{}{}

% ----------------------------------------------------------------
\begin{document}
\title[]
%{Convergence rate of numerical scheme finding extremal functions of Sobolev inequalities}
{test - 1}

\maketitle

\section{}

\subsection{}

Quadratic Programming (QP) problems are the following form:
\begin{align}
    \min_x \frac{1}{2} x^T Q x + q^T x \\
    \rm s.t. \quad & Gx \le g \\
    & Fx = f.
\end{align}
where $Q \in \R^{n \times n}$ positive semi-definite. $F \in \R ^{p \times n}$, $G \in \R ^{q \times n}$,
and $x \in \R^n$, $q \in \R^n$, $f \in \R ^p$, $g \in \R ^q$ are vectors.
Using slack variable, we can derive the following equation:
\begin{align}
    \min_x \frac{1}{2} x^T Q x + q^T x + I_{\R^+}(z) \\
    \rm s.t. \quad & Ax + Bz = c
\end{align}
where
\begin{equation}
    I_{\R^+}(z) = \begin{cases}
        0 & \quad z \ge 0 \\
        \infty & \quad \rm otherwise
    \end{cases}
\end{equation}
and
\begin{equation}
    A = \begin{bmatrix}
        G \\ F
    \end{bmatrix},
    \quad
    B = \begin{bmatrix}
        I_p \\ 0
    \end{bmatrix},
    \quad
    c = \begin{bmatrix}
        g \\ f
    \end{bmatrix}.
\end{equation}
$I_p$ is an $p \times p$ identity matrix.

\end{document}

